% Superseded 2026-09-25.
\begin{abstract}
Artificial-intelligence models are increasingly deployed at commercial energy sites to support forecasting and autonomous decision-making across a wide range of applications. However, their predictive performance degrades over time as operating conditions evolve, seasonal patterns shift, charging infrastructure expands, and user behaviour changes---a phenomenon known as concept drift. Existing strategies remain insufficient for real-world operation: periodic full-history retraining is computationally expensive as the dataset grows, rolling-window retraining discards the seasonal knowledge encoded in older data, and na\"ive fine-tuning on recent data leads to catastrophic forgetting of previously learned operating regimes. Although continual learning (CL) has received significant academic attention, its transition from offline experiments to continuously operating cyber-physical energy systems remains largely unexplored. In particular, existing studies rarely address the complete engineering chain required for autonomous adaptation, including cloud deployment, model lifecycle management, validation mechanisms, integration with real-time controllers, data reliability, and long-term operation in a real energy environment.

This paper presents an autonomous continual-learning architecture deployed as an operational prototype at a commercial energy site in Luxembourg, integrating AI forecasting, cloud-based model adaptation, and real-time operational control. The framework consists of: (i) two forecasting models---a PV LSTM and an EV CNN-LSTM---adapted nightly through an EWC-inspired proximal regulariser that uses a Memory Aware Synapses (MAS)-style output-sensitivity importance estimator; (ii) a tolerance-gated promotion mechanism with fail-safe retention that deploys a candidate only if it does not regress beyond a 5\% margin on a recent holdout, bounding---but not eliminating---the risk of deploying a worse model; and (iii) a direct coupling between the continually updated forecasts and a convex Model Predictive Control (MPC) scheduler that optimises electric-vehicle charging decisions under real ENTSO-E day-ahead electricity prices.

Beyond the algorithmic application of continual learning, this work demonstrates the complete pathway from initial model training and cloud deployment to integration with real-time control infrastructure. We present the EWC-inspired objective together with its Bayesian motivation, while making explicit that the deployed importance estimator is a MAS-style output-sensitivity proxy rather than an empirical Fisher estimator, and present the corresponding convex MPC formulation. Using real Leneda smart-meter data from Luxembourg, a representative nightly cycle reduces PV forecast MAE by 21.83\% and EV demand forecast MAE by 67.87\% relative to the static production models, driving a convex MPC loop against real day-ahead prices. A controlled $12$-cycle, five-seed sequential experiment ($240$ adaptation jobs) further isolates retention from plasticity: measured by backward transfer on a fixed winter reference window, bounded experience replay resists seasonal forgetting whereas the output-sensitivity regulariser does not improve retention over na\"ive fine-tuning, a paired advantage that is consistent across all five seeds ($11.1$~kW mid-stream, $p{=}0.002$). The paper also reports practical engineering challenges---framework incompatibilities, model-conversion (ONNX export) failures, data-quality degradation, and data-provenance verification---that are largely absent from the continual-learning literature. The primary contribution is therefore an operational architecture and real-site case study for safely integrating continual learning into an autonomous forecasting and MPC-based energy-management pipeline, rather than a novel CL algorithm.

\end{abstract}